\section{Introduction}\label{sec:introduction}

The local deformation theory of a polystable holomorphic bundle is
finite dimensional, whereas its Hermitian--Einstein metric is obtained
from a nonlinear elliptic equation. Near a bundle with additional
automorphisms, the linearized equation develops small eigenvalues.
Several of these eigenvalues can tend to zero at unrelated rates.
An absolute error estimate for the curvature or the operator then gives
no relative information about the weakest directions. The same issue
arises when comparing the harmonic metric on bundle moduli with a
finite-dimensional quotient metric.

We prove that canonical harmonic Kuranishi coordinates provide one
finite-dimensional model for all three comparisons near any fixed
locally free polystable compact-K\"ahler background. No formality
hypothesis is imposed. The complete convergent obstruction is retained,
and the estimates remain uniform when the stabilizer type changes.
The essential geometric observation is that the quadratic terms of both
the Kuranishi reconstruction and the complementary Hermitian gauge are
orthogonal to harmonic one-forms. This gives a quadratic estimate for
the ambient symplectic form and for the complementary quotient form.
The finite balancing theorem of \cite[Theorem~2.1]{wuebben14}, the
relative spectral reduction, and minimization over gauge directions
then apply on the same entire local family.

\subsection{The finite data}
Let $(X,\omega_0)$ be a compact connected K\"ahler manifold of complex
dimension $n$, and fix a locally free polystable bundle
\begin{equation}\label{eq:split}
 E_0=\bigoplus_{i=1}^{\ell}F_i\otimes W_i
\end{equation}
with pairwise nonisomorphic stable factors of equal slope. Put
$r_i=\operatorname{rank}F_i$, $m_i=\dim W_i$, and
$R=\sum_i r_im_i$. Fix a direct-sum Hermitian--Einstein metric $h_0$.
Let $\omega_t$ be a smooth finite-dimensional family of K\"ahler forms,
with $\omega_0$ its central value. We fix the determinant holomorphic
structure. Degrees and slopes are normalized by
\[
 \deg_\omega F=\int_Xc_1(F)\wedge\frac{\omega^{n-1}}{(n-1)!},
 \qquad \mu_\omega(F)=\frac{\deg_\omega F}{\operatorname{rank}F}.
\]
Write $d=\ddbar_0$ on $\End E_0$, and use the unnormalized
$(h_0,\omega_0)$ inner products. Denote the harmonic projections and
Green operators by $H_q$ and $G_q$.

The harmonic space $V=H^1(X,\End E_0)$ is the full Ext-quiver
representation of $G=\prod_i\mathrm{GL}(W_i)$; its compact form is
$K=\prod_i\mathrm U(W_i)$. Orthonormal harmonic bases in each
$\Ext^1(F_i,F_j)$ give the arrow coordinates $x_a$. Loops and all
multiplicities are included. Set
\[
 V_{\det}=\ker\bigl(\tr:V\longrightarrow H^1(X,\OO_X)\bigr).
\]
For small $x\in V$ the canonical harmonic recursion defines
\begin{equation}\label{eq:recursion-intro}
 a(x)=x-d^*G_2(a(x)\wedge a(x)),\qquad
 \kappa(x)=H_2(a(x)\wedge a(x)).
\end{equation}
Put $Z=\kappa^{-1}(0)\cap V_{\det}$ and
$E_x=(E_0,\ddbar_0+a(x))$ for $x\in Z$.
The full analytic structure of $Z$ is relevant to deformation theory.
For quotient metrics we use its underlying reduced analytic space.

The flat moment map and central target are
\begin{equation}\label{eq:moment}
 \begin{split}
 \mu_i(x)&=\sum_{h(a)=i}x_ax_a^*-\sum_{t(a)=i}x_a^*x_a,\\
 \zeta_i(t)&=2\pi\bigl(r_i\mu_{\omega_t}(E_0)-\deg_{\omega_t}F_i\bigr),
 \qquad \sum_i m_i\zeta_i(t)=0 .
 \end{split}
\end{equation}
Here $h(a)$ and $t(a)$ denote the head and tail of an arrow.
The ordinary multiplicity trace pairs the moment covector with
Hermitian matrices. The spectral mass has different rank weights:
\begin{equation}\label{eq:finite-spectrum}
 \begin{gathered}
 \g_0=\left\{(\xi_i)\in\bigoplus_i\End W_i:
                  \sum_i r_i\tr\xi_i=0\right\},\qquad
 C_x\xi=[\xi,x],\\
 q_x(\xi)=\|C_x\xi\|^2,\qquad
 \|\xi\|_0^2=\Vol(\omega_0)\sum_i r_i\tr(\xi_i^*\xi_i).
 \end{gathered}
\end{equation}
Let $\nu_1(x)\le\cdots\le\nu_N(x)$ be its eigenvalues, with complex
multiplicity and zeros included, where $N=\sum_i m_i^2-1$.

Locally near the origin consider
\[
 \Q_t=(Z\cap\mu^{-1}(\zeta(t)))/K .
\]
Its neighborhoods are images of small balanced neighborhoods; a bounded
ball is not declared invariant under $G$. Refine the reduced complex
quotient into smooth local pieces with locally constant stabilizer type
and smooth lifts, refining also at analytic singularities. Their existence
is explained in Section~\ref{subsec:pieces}. On a tangent represented by a
balanced lift $v$ define
\begin{equation}\label{eq:finite-metric}
 g_{\mathrm{fin},x}([v],[v])
       =\inf_{\xi\in\g_0}\|v-C_x\xi\|^2.
\end{equation}
This is the flat K\"ahler quotient metric restricted to the chosen piece.
The tangents in this formula are genuine tangents to those pieces, not
arbitrary Zariski tangents of a singular coarse space.

\subsection{The local comparison theorem}
Let $b_t$ be the $R$th root of the ratio of the normalized
$\omega_t$-Hermitian--Einstein determinant metric to $\det h_0$.
We normalize it by $\int_X\log b_t\,dV_0=0$, where
$dV_0=\omega_0^n/n!$.
Then $b_0=1$ and $\|\log b_t\|_{C^j}\le C_j|t|$.

\begin{theorem}\label{thm:main}
For the fixed data above, there are a parameter neighborhood, nested
coefficient neighborhoods, and an open neighborhood of the central
actual Hermitian--Yang--Mills class with the following properties.
Write $\delta=|t|+\|x\|^2$ and shrink the neighborhoods so that
$C\delta<1/2$.
\begin{enumerate}
\item Every $x\in Z\cap\mu^{-1}(\zeta(t))$ in the inner neighborhood
has a Hermitian--Einstein metric with the prescribed normalized
determinant and
\begin{equation}\label{eq:metric-main}
 e^{-C\delta}b_th_0\le h^{\mathrm{HE}}_{t,x}
                         \le e^{C\delta}b_th_0.
\end{equation}
The construction identifies neighborhoods in $\Q_t$ and the actual
fixed-determinant HYM quotient, uniformly for small $t$, and covers
every sufficiently nearby actual solution. Metrics are unique modulo
determinant-preserving holomorphic automorphisms; the unitary HE class
is unique.

\item The Dolbeault Laplacian on $\End_0E_x$, for
$h^{\mathrm{HE}}_{t,x}$ and $\omega_t$, has exactly $N$ eigenvalues
below a fixed $\gamma>0$, and
\begin{equation}\label{eq:spectrum-main}
 (1-C\delta)\nu_j(x)\le\lambda_j(t,x)\le
 (1+C\delta)\nu_j(x)\quad(1\le j\le N).
\end{equation}
Its kernel is naturally identified with $\ker C_x$, and its next
eigenvalue is at least $\gamma$. For full $\End E_x$, add the one
exact scalar zero; all other scalar modes have a uniform positive gap.

\item The local identification $f_t$ is a homeomorphism, preserves the
local complex-orbit description, and is smooth on the smooth pieces
specified above. Their harmonic Kodaira--Spencer metric satisfies
\begin{equation}\label{eq:quotient-main}
 (1-C\delta)g_{\mathrm{fin}}\le f_t^*g_t^{L^2}
                                      \le(1+C\delta)g_{\mathrm{fin}} .
\end{equation}
The same comparison constant works through changes of stabilizer.
\end{enumerate}
Constants and radii may depend on the fixed background, ranks,
multiplicities, harmonic dimensions, positive background elliptic
gaps, Sobolev and elliptic constants, the selected regularity order,
and bounds for the given family $\omega_t$. They are independent of
every positive $\nu_j(x)$, ratios of positive finite eigenvalues,
positive singular values of $C_x$ or $D\kappa_x$, and stabilizer type.
\end{theorem}

For $N=0$ and $R>1$, the trace-free assertion is a uniform gap on the
entire zero-form spectrum. For $R=1$, $\End_0E_0=0$ and that assertion
is vacuous; the full endomorphism statement retains its scalar zero.
The theorem asserts no uniform radius for separate charts on smooth
pieces, nor a continuous tensor on untransported tangent vectors when
their dimension changes.

The error can be sharpened when the harmonic coordinates move with the
polarization. Continue the HE metric of each fixed stable factor
smoothly to $h_i(t)$ and use the harmonic construction for
$h_{\mathrm{gr},t}=\bigoplus_i h_i(t)\otimes I_{W_i}$ and $\omega_t$.
Proposition~\ref{prop:moving} proves all three comparisons with error
$C\|x\|_t^2$, relative to that background and its spectral mass.
The split bundle need not remain polystable for $t\ne0$.
This distinction preserves the stronger rates of the earlier extension
theorems when a polarization moves predominantly along a stability wall.

\subsection{An application with cubic obstruction}
We construct a stable rank-two bundle $A$, with trivial determinant,
on a smooth bicubic threefold in $\PP^2\times\PP^2$, whose analytic
deformation germ is
\begin{equation}\label{eq:cubic-intro}
 \Specan\C\{a,b\}/(a^2b,ab^2).
\end{equation}
Its quadratic obstruction vanishes, but mixed directions are obstructed
at order three. Section~\ref{sec:bicubic} proves smoothness of the compact
geometry, stability, and equality of the complete curve and bundle
deformation germs. The construction modifies the bicubic Serre method
of Thomas \cite{thomas}; its deformation equations differ from his
corrected quadratic example.

Theorem~\ref{thm:main} applies to the entire $18$-dimensional harmonic
deformation problem at $A\otimes\C^3$, retaining its full matrix-valued
obstruction. The trace-free finite algebra has dimension eight.
An explicit diagonal subfamily has two eigenvalues on an arbitrarily
weak positive scale and four on a larger scale, with changing exact
kernel dimensions. This subfamily illustrates spectral uniformity;
the application concerns the whole local moduli problem.
The reduced space of \eqref{eq:cubic-intro} is the node $ab=0$.
The harmonic metric is defined on reduced smooth pieces and does not
measure its embedded nilpotent structure.

\subsection{Relation to earlier work and proof strategy}
The Donaldson--Uhlenbeck--Yau correspondence gives the qualitative
existence of HE metrics \cite{donaldson,uy}; see also \cite{kobayashi}.
Buchdahl--Schumacher develop the equivariant Kuranishi and automorphism
description \cite{bs22}. Their Proposition~6.2 and Lemma~6.3 already
give a quadratic complementary correction and a quartic finite moment
remainder. Their Theorems~7.2 and~8.1 provide quantitative raywise
correction and common existence neighborhoods. The present comparison
requires relative control in each weak orbit direction, which does
not follow from an absolute quartic remainder.
Their treatment of harmonic moduli metrics also includes obstructed
spaces under the hypotheses of \cite[Theorem~11]{bs21}.
Delloque's local correspondence supplies the common-parameter slices,
morphism description and quotient topology used here
\cite[Section~3.2, Propositions~4.3, 4.5, 5.6,
Corollary~5.2 and Theorem~7.4]{delloque}.
Convergent Ext-quiver relations are also part of Toda's local
description of sheaf moduli \cite[Theorem~1.1]{toda}.

Several quantitative mechanisms are available in earlier special
geometries. The complete multiplicity algebra and relative spectral
reduction occur in \cite{wuebben4}; complementary minimization and
nonabelian quotient estimates occur in \cite{wuebben6,wuebben7}.
The finite balancing theorem of \cite{wuebben14} already applies to
arbitrary quivers and locally invariant analytic subsets.
The additional result here is the general geometric verification
that permits these mechanisms to act simultaneously on the entire
convergent Kuranishi family, without formality or a positive scale
ratio bound. The cubic compact example provides a locally free
application beyond the earlier extension and K3 settings.

The proof begins with the two orthogonal quadratic jets
\[
 -d^*G_2(x\wedge x),\qquad -d\sigma_2(x).
\]
Their pairing with harmonic forms vanishes, giving the quadratic
ambient estimate. Finite balancing then solves the remaining moment
equation with a small parallel gauge. For spectra, the distinct identity
$D_a\xi=-D_xa(C_x\xi)$ controls the off-diagonal Schur term by the
finite energy of the same vector. For quotient metrics, minimizing
first over the uniformly coercive infinite-dimensional complement
gives a positive form close to the harmonic form; only then is the
finite orbit minimum taken. Derivatives of the balancing gauge are
vertical and disappear. These separate arguments explain why one
relative constant survives arbitrarily weak finite directions.

